\chapter{Pourquoi une démarche}\label{ch:demarche}

\begin{objectifs}[Objectifs]
\begin{itemize}
  \item expliquer en trois phrases pourquoi un projet logiciel dérape ;
  \item citer les six principes de la démarche et leur garde-fou ;
  \item situer chacune des sept étapes et son skill ;
  \item dire ce que le plugin ne fera jamais à votre place.
\end{itemize}
\end{objectifs}

\begin{situation}[Sur le terrain : \sika{}]
Un mardi, la présidente de l'Association Entraide Dantokpa reçoit une équipe de développeurs.
Elle décrit son problème en une phrase : \q{Chaque semaine, quelqu'un dit qu'il a payé et le cahier
dit le contraire.} Dix minutes plus tard, un développeur propose déjà un écran de connexion, une
messagerie interne et un classement des meilleures cotisantes. Personne n'a demandé combien de
cotisations sont contestées aujourd'hui. Ce chapitre explique comment éviter cette scène.
\end{situation}

\section{Le problème}
Un projet logiciel dérape rarement à cause du code. Il dérape parce qu'on construit la mauvaise
chose, parce qu'on se comprend mal, ou parce qu'on décide trop tôt sans le dire. La démarche
Lumière répond à ces trois risques par une chaîne d'étapes courtes ; chacune produit un
livrable \emph{validé par un humain} avant que la suivante ne commence.

Le plugin \textbf{Simetriik Forge} met cette démarche en pratique avec un assistant de code. Il
apporte deux choses, à ne pas confondre :
\begin{description}
  \item[Des skills] des instructions que l'assistant suit pour une étape (poser les bonnes questions, remplir le bon gabarit). Ils \emph{guident}.
  \item[Des garde-fous] de petits scripts (\cmd{scripts/}) qui \emph{échouent} quand une règle est violée. Ils \emph{contrôlent}, sans interprétation, et fonctionnent aussi en intégration continue.
\end{description}

\begin{concept}[Garde-fou]
Un script déterministe qui vérifie une règle de la démarche et retourne un code d'échec si elle
est violée. Il ne juge pas la qualité : il constate un fait (un statut, un nom, un fichier
présent). La qualité reste l'affaire de la relecture humaine.
\end{concept}

\section{Les six principes}
\noindent\begin{tabularx}{\linewidth}{>{\raggedright\arraybackslash}p{6.3cm}X}
\toprule
\textbf{Principe} & \textbf{Ce qui le fait respecter} \\ \midrule
Pas de code sans spécification validée. La validation est humaine. & \cmd{check-spec-validated.sh}, hook \emph{pre-commit} \\
Chaque fonctionnalité remonte à un objectif mesurable. & l'impact map, colonne « Objectif » obligatoire \\
Un terme = une définition = un nom technique. & \cmd{check-glossary-terms.sh} (FF-05) \\
Le contrat d'API précède le code. & \cmd{check-openapi-first.sh} \\
Les décisions structurantes sont consignées, stack comprise. & \cmd{check-adr-present.sh}, ADR \\
Ce qui s'automatise s'automatise. & les scripts, en local et en CI \\
\bottomrule
\end{tabularx}

\section{Les sept étapes d'un coup d'œil}
La démarche compte sept étapes numérotées de 0 à 6. La mise en œuvre du code, quant à elle,
se glisse entre le contrat et la recette.

\begin{figure}[H]
\centering
\resizebox{\linewidth}{!}{%
\begin{tikzpicture}[
  st/.style={draw=accent,fill=accent!7,rounded corners=3pt,minimum width=3.1cm,minimum height=15mm,align=center,font=\sffamily\small,text=ink},
  co/.style={draw=ochre,fill=ochre!10,rounded corners=3pt,minimum width=3.1cm,minimum height=15mm,align=center,font=\sffamily\small,text=ink},
  num/.style={circle,fill=accent,text=white,font=\sffamily\bfseries\footnotesize,inner sep=2pt},
  ar/.style={-{Stealth[length=2.2mm]},thick,accent}]
  \node[st] (a) at (0,0)   {\textbf{Cadrage}\\[-1pt]{\scriptsize forge-cadrage}};
  \node[st] (b) at (3.8,0) {\textbf{Domaine}\\[-1pt]{\scriptsize forge-domaine}};
  \node[st] (c) at (7.6,0) {\textbf{Spécification}\\[-1pt]{\scriptsize forge-spec}};
  \node[st] (d) at (11.4,0){\textbf{Architecture}\\[-1pt]{\scriptsize forge-architecture}};
  \node[st] (e) at (11.4,-2.6){\textbf{Contrat d'API}\\[-1pt]{\scriptsize forge-contrat}};
  \node[co] (f) at (7.6,-2.6){\textbf{Code}\\[-1pt]{\scriptsize forge-implemente}};
  \node[st] (g) at (3.8,-2.6){\textbf{Recette}\\[-1pt]{\scriptsize forge-recette}};
  \node[st] (h) at (0,-2.6){\textbf{Production}\\[-1pt]{\scriptsize forge-prod}};
  \node[num] at (a.north west) {0}; \node[num] at (b.north west) {1}; \node[num] at (c.north west) {2};
  \node[num] at (d.north west) {3}; \node[num] at (e.north west) {4}; \node[num] at (g.north west) {5}; \node[num] at (h.north west) {6};
  \draw[ar] (a)--(b); \draw[ar] (b)--(c); \draw[ar] (c)--(d); \draw[ar] (d)--(e);
  \draw[ar] (e)--(f); \draw[ar] (f)--(g); \draw[ar] (g)--(h);
  \draw[ar,dashed,grey] (h.south) -- ++(0,-8mm) -| node[pos=.25,below,font=\sffamily\scriptsize,text=grey]{lot suivant : retour au cadrage} ([xshift=-8mm]a.south west) |- (a.west);
\end{tikzpicture}}
\caption{Le chemin de la démarche : sept étapes validées, et le code entre le contrat et la recette.}
\end{figure}

À côté de ce chemin, trois skills transversaux : \skill{forge-init} (démarrer), \skill{forge-statut}
(où en suis-je ?) et \skill{forge-rapide} (le cycle court des petites fonctionnalités, chapitre~\ref{ch:rapide}).

\section{Ce que le plugin ne fait pas}
Il ne déploie jamais en production de lui-même, ne valide aucune spécification, ne pousse aucun
code sans demande et n'invente aucune règle métier : quand une information manque, il pose la
question. C'est volontaire. La validation engage une responsabilité (celle du client, ou la
vôtre) ; elle ne se délègue pas à un outil.

\begin{piege}[Piège : « l'assistant l'a validé »]
Un assistant peut passer un statut à \emph{Validée} en une ligne ; le script, lui, le croira.
C'est pourquoi le plugin lui interdit de le faire : la ligne \emph{Validé par} porte un nom et
une date \emph{humains}. Une spec validée par personne est une spec non validée.
\end{piege}

\begin{vousjouez}[À vous de jouer]
Prenez un projet que vous avez vécu et qui a dérapé. En une phrase chacune, laquelle de ces trois causes
l'explique le mieux : \emph{mauvaise chose construite}, \emph{malentendu}, \emph{décision tardive non dite} ?
Notez-la dans votre carnet : vous y reviendrez aux chapitres~\ref{ch:cadrage}, \ref{ch:domaine} et \ref{ch:archi}.
\end{vousjouez}

\begin{retenir}[À retenir]
\begin{itemize}
  \item Trois risques : mauvaise chose, malentendu, décision tardive.
  \item Les \emph{skills} guident ; les \emph{garde-fous} contrôlent ; l'\emph{humain} valide.
  \item Sept étapes, un livrable validé chacune, le code entre le contrat et la recette.
\end{itemize}
\end{retenir}
\source{README.md}

\begin{acquis}
  \item Quelle différence y a-t-il entre un skill et un garde-fou ?
  \item Qui peut passer une spécification au statut \emph{Validée} ?
  \item Dans quel ordre viennent le contrat d'API et le code ?
\end{acquis}
